% from ega4-18-fr.tex
\oldpage[IV-4]{129}

Reste à montrer que sous l'hypothèse $a)$, l'application (18.5.4.2) est \emph{surjective} ($S$ étant quasi-compact et quasi-séparé). Soit donc $u_0:S_0\to S'_0$ une $S_0$-section de $S'_0$; $u_0(S_0)$ étant quasi-compact dans $S'$, peut être recouvert par un nombre fini d'ouverts affines $V_i$ tels que la restriction $g|V_i$ soit un morphisme de type fini; si $V$ est la réunion des $V_i$, on en conclut que $g|V:V\to S$ est un morphisme \emph{de présentation finie}, étant séparé et localement de présentation finie par hypothèse (1.6.1). Remplaçant $S'$ par $V$, on peut donc supposer que $g$ est de présentation finie. Comme $S$ est quasi-compact et quasi-séparé et $g$ quasi-fini et séparé (17.6.1), il résulte du « Main theorem » (8.12.6) que $g$ se factorise en $S'\xrightarrow{j}S''\xrightarrow{f}S$, où $j$ est une immersion ouverte et $f$ un morphisme \emph{fini}. Posons $S''_0=S''\times_S S_0$, $j_0=j_{(S_0)}:S'_0\to S''_0$, qui est une immersion ouverte et $f_0=f_{(S_0)}:S''_0\to S_0$, qui est un morphisme fini. Alors $u_0$ est aussi une $S_0$-section de $S''_0$. Comme $g_0:S'_0\to S_0$ est étale, $u_0$ est une immersion ouverte de $S_0$ dans $S'_0$ (17.4.1), donc $u_0(S_0)$ est ouvert dans $S'_0$, et \emph{a fortiori} dans $S''_0$; mais d'autre part, comme $f_0$ est un morphisme fini, donc séparé, $u_0$ est une immersion fermée de $S_0$ dans $S''_0$ (I, 5.4.6), donc $X_0=u_0(S_0)$ est \emph{à la fois ouvert et fermé} dans $S''_0$. En vertu de l'hypothèse $a)$, il existe un sous-ensemble à la fois ouvert et fermé $X$ de $S''$ tel que $X\cap S''_0=X_0$. Montrons d'abord que le morphisme $f:S''\to S$ est \emph{étale aux points de $X$}: en effet, l'ensemble $U$ des points de $X$ où $f|X$ est étale est ouvert et contient par hypothèse $X_0\subset S'_0$. Mais le lemme (18.5.4.3) appliqué au morphisme fini $f|X$ prouve que $U=X$. D'autre part, $S'\cap X$ est ouvert dans $X$ et contient $X_0$ par hypothèse, donc le même raisonnement prouve que $S'\cap X=X$, c'est-à-dire $X\subset S'$. Il reste à montrer que, pour tout $s\in S$, le nombre géométrique $n(s)$ de points de $X\cap f^{-1}(s)$ est égal à $1$, car il en résultera que $f|X$ est radiciel et surjectif, et comme $f|X$ est étale, on aura montré (17.9.1) que $f|X$ est un isomorphisme de l'ouvert $X\subset S'$ sur $S$, dont l'isomorphisme réciproque $u$ sera la $S$-section cherchée prolongeant $u_0$. Mais comme $f|X$ est étale et \emph{fini}, $s\mapsto n(s)$ est continue dans $S$ (18.2.8), et comme $X\cap S''_0=X_0$, on a $n(s)=1$ dans $S_0$; l'ensemble des points $s\in S$ tels que $n(s)=1$ étant ouvert dans $S$ et contenant $S_0$, il est égal à $S$ par (18.5.4.3). C.Q.F.D.

\begin{remark}[18.5.4.4]
\label{IV.18.5.4.4}
On peut montrer que l'énoncé (18.5.4) reste valable lorsque, dans la condition $b)$, on suppose seulement le morphisme $g$ étale (mais non nécessairement séparé) [43, exp. XII].
\end{remark}

\begin{definition}[18.5.5]
\EGAReaderTarget{ega4.c04.def.18-5-5}{IV.18.5.5}
On dit qu'un préschéma $S$ et un sous-préschéma fermé $S_0$ de $S$ forment un \emph{couple hensélien} s'ils vérifient la condition $a)$ de (18.5.4).
\end{definition}

Compte tenu de (I, 5.1.8), il revient au même de dire que $(S,S_0)$ est un couple hensélien ou que $(S_{\mathrm{red}},(S_0)_{\mathrm{red}})$ en est un.

\begin{proposition}[18.5.6]
\label{IV.18.5.6}
\begin{enumerate}
  \item[(i)] Si $(S,S_0)$ est un couple hensélien, alors, pour tout morphisme fini $f:S'\to S$, si $S'_0$ est le sous-préschéma $f^{-1}(S_0)$ de $S'$, le couple $(S',S'_0)$ est hensélien.
  \item[(ii)] Soient $S=\coprod_\alpha S^{(\alpha)}$ une somme de préschémas, $S_0$ un sous-préschéma fermé de $S$, somme des sous-préschémas fermés $S_0^{(\alpha)}$ des $S^{(\alpha)}$. Pour que le couple $(S,S_0)$ soit hensélien, il faut et il suffit que chacun des couples $(S^{(\alpha)},S_0^{(\alpha)})$ le soit.
\end{enumerate}
\end{proposition}

L'assertion (i) est conséquence immédiate de la définition, puisque pour tout morphisme fini $g:S''\to S'$, $f\circ g:S''\to S$ est un morphisme fini. De même, sous les conditions de (ii), pour qu'un morphisme $g:S'\to S$ soit fini, il faut et il suffit que chacune

